\chapter{Sanitizers}
In questo capitolo esploreremo le vulnerabilità critiche legate alla gestione 
della memoria in C/C++ (come \textit{buffer overflow}, \textit{use-after-free} e \textit{uninitialized 
memory}) e analizzeremo il funzionamento dei "Sanitizers". In 
particolare, metteremo in luce la contrapposione tra l'approccio dinamico di \textbf{Valgrind} a quello statico 
di \textbf{AddressSanitizer} (\textbf{ASan}), spiegandone l'implementazione a livello di 
shadow memory e l'impatto sulle performance.

I bug di gestione della memoria sono subdoli perché un programma può 
tecnicamente "sopravvivere" a un buffer overflow senza andare in crash.
Questo accade per ragioni architetturali specifiche:

\begin{itemize}
    \item La memoria heap è pre-allocata ed espandibile, con permessi di lettura e scrittura (R/W).
    \item La memoria viene allocata dal sistema operativo in "\textbf{pagine}" (tipicamente da 4096 byte su architetture Linux/x86).
    \item Se un overflow eccede il buffer allocato (es. 30 byte) ma rimane 
    all'interno della stessa pagina di memoria, non viene generata alcuna 
    eccezione a livello di CPU (\textit{page fault}) e il programma continua 
    l'esecuzione.
    \item Il crash si verifica solo quando l'overflow tenta di accedere a una pagina non mappata o protetta, causando un'eccezione hardware.
\end{itemize}

\section{Analisi delle vulnerabilità legate alla memoria}

\subsection{Buffer Over-read}

A differenza dell'overflow distruttivo, l'over-read non sovrascrive nulla, 
ma legge oltre i limiti previsti, causando la fuga di dati sensibili.
Esempio di buffer Over-read storicamente noto è il bug \textbf{Heartbleed} di 
\textit{OpenSSL} nel quale il client invia un payload contenente una stringa
e la sua lunghezza dichiarata; se il client invia ad esempio la parola "HAT" 
ma dichiara una lunghezza di 500 caratteri, il server alloca un buffer e 
copia i 3 caratteri. Successivamente, il server legge 500 byte a partire 
dall'indirizzo di memoria di "HAT", inviando indietro al client i 497 byte 
adiacenti che si trovavano già in memoria; il problema è proprio legato a quei
497 byte poiché possono contenere chiavi master del server, password di altri 
utenti o dati sensibili.

\subsection{Double-free \& Use-After-Free}
Come sappiamo in C/C++ non abbiamo un \textit{garbage collector} infatti gestiamo
manualmente la memoria. Non dovremmo assolutamente fare due volte una \textit{free} e
nè dovremmo usare un puntatore ad una certa area se abbiamo fatto la \textit{free}.

Questa vulnerabilità nasce proprio quando la memoria allocata dinamicamente 
(con \textit{new} o \textit{malloc}) viene liberata (con \textit{free} o 
\textit{delete}) ma il programma continua ad utilizzare un puntatore che punta 
a quell'area.

Esempio di \textbf{Safari} (Apple):

\begin{lstlisting}
//un oggetto "animator" viene creato
Animator * animator = getAnimatorObj();
...
//la funzione setCurrentTime() attiva l'handler() che resetta l'oggetto "animator"
...
delete svg->animator_ptr;
//...continuo della funzione setCurrentTime()
animator->DoSomething(); // uso della memoria liberata!
\end{lstlisting}

Nel motore JavaScript, viene creato un oggetto \textit{animator\_ptr} in C++. 
Una funzione \textit{setCurrentTime()} innesca una reazione a catena che 
cancella (\textit{delete}) questo oggetto, tuttavia, il puntatore viene 
riutilizzato subito dopo chiamando \textit{animator->DoSomething()}.

L'attaccante a questo punto può iniettare uno script JS \textit{MaliciousObject} poiché 
l'allocatore heap dell'OS, per efficienza, riutilizza l'area di memoria 
appena liberata assegnandola al \textit{MaliciousObject}. Quando viene 
chiamato \textit{animator->DoSomething()}, il programma esegue in realtà 
il codice dell'attaccante.

\subsection{Uninitialized memory}

Le variabili non inizializzate in C contengono i "rifiuti" (dati precedenti) 
rimasti in quel segmento di memoria.

Esempio del \textbf{kernel Linux} (CVE-2016-4569): tale problema dimostra 
come le chiamate di sistema (\textit{system call}) spostino lo stack 
pointer del kernel, ma non puliscano i contenuti dello stack. Se il kernel 
non inizializza una struttura dati e la copia in user space, i byte non 
inizializzati conterranno frammenti del kernel stack, permettendo agli 
attaccanti di bypassare mitigazioni come \textbf{ASLR}.

\subsection{Format string attacks}

Le funzioni della famiglia printf() e scanf() utilizzano le \textbf{format string}, ovvero un
linguaggio minimale per indicare e definire input e output, tipo \%s, \%d e così via.
Il problema nasce dal fatto che alcuni programmi permettono all'attaccante 
di controllare la format string, vediamo un esempio:

\begin{lstlisting}
int main() {
    char nome[100];
    printf("Inserisci un nome: ");
    fgets(nome, sizeof(nome)-1, stdin);
    printf(nome);       //SBAGLIATO!
    printf("%s", nome); //OK!
}
\end{lstlisting}

Dunque per fare la stampa di una variabile va usata sempre la format string "\%s" nella \textit{printf()}. 
Se non si fa e si fa direttamente \textit{printf(nome)}, l'attaccante, nella stringa \textit{nome}, 
può mette valori speciali come "\%x\%x\%x\%x" che gli consente di stampare lo
stack oppure "\%n" per sovrascrivere un'area di memoria arbitraria.

\subsection{Integer overflow}

Gli interi in C usano un numero fisso di bit; se si supera il limite massimo, 
il valore "wrappa" a un numero negativo o tronco.

\begin{lstlisting}
int check(unsigned short len, char * ptr) {
    if(len >= 10) return -1;
    printf("s = %d\n", len);
    return 0;
}

int main(int argc, char *argv[]) {
    char buf[10];
    int n = atoi(argv[1]);
    if( check(n, argv[2]) == -1 ) {
        printf("Oh no you don't!\n");
        return -1;
    }
    strncpy(buf, argv[2], n);
    printf("%s\n", buf);
    return 0;
}
\end{lstlisting}

Nell'esempio riportato notiamo che nella funzione \textit{check()} abbiamo 
un \textit{implicit cast} ad \textbf{unsigned short} ed, essendo 65535 il massimo 
valore per uno \textit{short int}, avremo un integer overflow nel caso in cui n sia maggiore
di quel valore.

Spesso attacchi di questo tipo vengono utilizzati per arrivare ad un buffer
overflow, ad esempio attaccando l'indice di un array per farlo eccedere la sua
dimensione massima e quindi accedere a più dati.

\section{Sanitizer}

Sono tool che hanno il ruolo di bug detector per C/C++, sono capaci cioè di darci un
avviso nel momento in cui il nostro programma ha un problema di quelli visti in
precedenza nella gestione della memoria. Il programma o viene fermato o 
comunque veniamo avvisati del tutto.

Fondamentalmente si va ad inserire delle sonde (dunque codice aggiuntivo) nel nostro
programma, definite \textbf{instrumentation}, utilizzate per capire se ci sono dei problemi
nell'utilizzo della memoria. Il programma originale viene modificato 
aggiungendo le sonde per collezionare informazioni riguardo l'esecuzione e 
per prendere decisioni sulla base di queste.

\begin{itemize}
    \item Se la code instrumentation va a prendere il codice macchina prima che venga
    eseguito sul processore e inserisce delle sonde si parla di \textbf{instrumentazione
    dinamica}, in quanto queste vengono iniettate a run-time.
    \item Se l'instrumentazione avviene al tempo di compilazione, ovvero le sonde
    vengono inserite prima della compilazione nel codice sorgente, si parla allora
    di \textbf{instrumentazione statica}.
\end{itemize}
Esempio di instrumentazione statica:

\begin{lstlisting}
int main() {
    allocations.init();
    char * p = malloc(100);
    allocations.insert(p, 100);
    int offset = 200;
    if(allocations.find(p+offset)) print_error(p+offset);
    p[offset] = 'A';
    allocations.find_leaks();
}
\end{lstlisting}

Nell'esempio ogni volta che si crea un oggetto viene aggiunta un'istruzione che si
segna l'oggetto creato e la sua dimensione. Nel momento dell'utilizzo si aggiunge
prima di questo un \textit{if} che verifica che l'offset inserito non ecceda la dimensione
massima segnata. Prima della fine del programma si fa anche una 
ricerca dei \textit{memory leak}, ovvero si controlla che tulle le variabili allocate 
siano state correttamente deallocate.

Anche la stessa \textbf{Code Coverage} viene fatta con lo stesso principio, inserendo
delle sonde e valutando alla fine delle esecuzioni quante di queste sono state
eseguite.

\subsection{Valgrind (Memcheck)}

Valgrind è un framework di analisi dinamica che non richiede la 
ricompilazione del codice sorgente Esso utilizza la traduzione \textbf{Just-In-Time} 
(JIT) per eseguire il binario all'interno di una CPU virtuale, iniettando 
istruzioni di analisi a runtime. \textbf{Shadow Memory} (\textbf{Memcheck}): mantiene una 
mappa della memoria usando "\textit{V bits}" (\textit{Valid-value}, indica se il dato è 
inizializzato) e "\textit{A bits}" (\textit{Valid-address}, indica se l'indirizzo è allocato 
legittimamente). Il bit V segue il dato ovunque, persino dentro i registri 
della CPU virtuale.

\textbf{Limiti di Valgrind}:
\begin{itemize}
    \item \textbf{Overhead estremo}: Valgrind rallenta l'esecuzione da 10x a 60x e triplica l'uso di RAM.
    \item \textbf{Blind spot sugli array di stack}: Avendo ad esempio due buffer \textit{int a[16]; int b[16];} sullo 
    stack, essi sono contigui in memoria. A livello binario (assembly), accedere 
    a \textit{a[16]} significa calcolare un offset dal registro \textit{rbp}. Valgrind vede 
    solo un accesso legittimo alla memoria dello stack e non riesce a 
    distinguere i confini tra le due variabili, fallendo nel rilevare 
    l'overflow locale.
\end{itemize}

\subsection{AddressSanitizer (ASan)}

ASan risolve molti dei limiti di Valgrind operando a compile-time come 
passaggio aggiuntivo del compilatore tramite il flag \textit{-fsanitize=address}.

Dato che ASan opera durante la compilazione, altera il layout della memoria 
del programma; aggiunge zone "cuscinetto" (\textit{redzones}) attorno a ogni variabile 
globale, buffer nell'heap o array nello stack (tipicamente di 32 byte). 
Queste aree vengono marcate come inaccessibili nella \textit{shadow memory}. A 
proposito della \textit{shadow memory} va detto che ASan mappa ogni word 
di 8 byte del programma in un signolo byte e dunque fa un uso più efficiente 
della memoria rispetto a Valgrind, inoltre l'overhead sulla CPU è mediamente 
solo 2x, rendendolo utilizzabile anche in contesti di test estesi.

In generale, il compilatore inietta un controllo primordiale prima di ogni 
istruzione di lettura o scrittura:

\begin{lstlisting}
// Codice iniettato da ASan
Shadow = (Addr >> 3) + Offset; //Shadow byte
if (*Shadow != 0) ReportAndCrash(Addr); //Crasha se lo stato e' "not addressable"
// Codice originale
*Addr = ...
\end{lstlisting}

Se un overflow cerca di toccare la \textit{redzone}, il controllo sul byte 
ombra (\textit{shadow byte}) fallisce immediatamente e il programma va in crash volontario 
riportando il bug.

ASan può riconoscere anche il problema dello \textit{use-after-free} in quanto
ad ogni variabile liberata il tool marca quei byte come "\textit{quarantined}"; quindi
se si prova successivamente ad accedere a quella zona ASan segnala che si sta
accedendo ad una zona messa in quarantena.

\textbf{Limiti di ASan}
\begin{itemize}
    \item Più leggero, ma \textbf{meno preciso} di Valgrind avendo una \textbf{granularità di 8 byte}.
    \item Nessuna distinzione tra i permessi di lettura e scrittura.
    \item ASan può controllare solo gli accessi allineati in memoria.
\end{itemize}

In conclusione possiamo dire che entrambe le suite di tool possono rilevare
la maggior parte dei bug relativi alla gestione della memoria, in particolare:

\begin{itemize}
    \item ASan offre un miglior supporto per il rilevamento dello \textit{stack/global array overrun}.
    \item Valgrind/Memcheck offre un miglior supporto per il rilevamento dell'\textit{uninitialized memory}.
\end{itemize}